\emph{Setting.} The gravitational field is $\hb(b) = b/a(\eta)$ on flat spacetime (Definition \claimref{obj.synthesis.gtg-flat-spacetime-gauge-fields}), $\eta = x\cdot\gamma_{0}$, with $a = a'\eta$, $a' > 0$ constant: the radiation-type descent, whose wall $\eta = 0$ is the radical-normal tame wall in $t = \eta^{2}/4$ (Theorem \claimref{synthesis.radical-normal-tame-wall-is-a-regular-initial-surface-whose-smoothness-selects-no-state}(iv)). Hypothesis (DF) continues the chart to $\eta < 0$ with the same formula, so that the rescaled field $\chi$ ($\chi = a^{3/2}\psi$ on $\eta > 0$) obeys $\nabla\chi\,I\sigma_{3} = M\chi\gamma_{0}$, $M = \mu\eta$, $\mu = ma' \ge 0$, on both sides. For a plane wave $\chi = \phi(\eta)\exp(-I\sigma_{3}\,\vec k\cdot\vec x)$, $\phi$ even, the cited result's computation gives $\phi' = G\phi$, $G = MA_{0} + B_{k}$, $A_{0}\phi = \gamma_{0}\phi\gamma_{0}(I\sigma_{3})^{-1}$, $B_{k}\phi = -\gamma_{0}\kappa\phi(I\sigma_{3})^{-1}$, $\kappa = k_{i}\gamma_{i}$, with $G^{2} = -(k^{2} + M^{2})$; the check recomputes this on the eight even blades at four samples and verifies that $G$ is antisymmetric for the scalar product $\langle\gamma_{0}\tilde\phi\gamma_{0}\chi\rangle_{0}$ (the Gram matrix of the even blades is the identity), so the norm is conserved. Each momentum and helicity sector is the stored two-level system $Iu' = Hu$, $H = M\sigma_{3} + k\sigma_{1}$, in the Pauli algebra $\Cl(3,0)$ with unit its pseudoscalar $I$, and for $\rho = \tfrac12(1 + s)$ it is $s' = -2I(h\wedge s)$, $h = k\sigma_{1} + M\sigma_{3}$ (Theorem \claimref{synthesis.massive-dirac-instantaneous-vacuum-and-non-smooth-descents-are-not-hadamard}, Setting): the mass part of $G$ is the $\sigma_{3}$ part of $H$ and the momentum part the $\sigma_{1}$ part. The equation is linear and its flow is a rotation, since $(v - r)\cdot\bigl(-2I(h\wedge(v - r))\bigr) = 0$: scalar products of solutions are constant in $\eta$ (the check verifies it coefficientwise on power series to order $24$).

\emph{The class of states.} A gauge-invariant quasi-free state is fixed by its one-particle operator $0 \le \Gamma \le 1$ on the solution space. Homogeneity makes $\Gamma$ commute with the spatial translations; the class considered in (i) to (vi) is isotropic, diagonal in helicity with the same vector in the two helicities, and charge-symmetric (equal occupation of particles and antiparticles), so on each sector $\Gamma = \tfrac12(1 + v_{k})$ with $v_{k}$ a vector of the Pauli algebra depending on $k = |\vec k|$ only, $|v_{k}| \le 1$, solving the Bloch equation; in (vii) the vector may depend on the sector; the state is a Fock (pure) state exactly when $|v| = 1$ on every sector. The expected number of quasi-particles of a pure reference $r$ in the sector is $\tfrac12(1 - v\cdot r)$, the weight of $\tfrac12(1 + v)$ on the complement $\tfrac12(1 - r)$ (the stored formula $|\beta|^{2} = \tfrac12(1 - s\cdot s')$, extended linearly to $|v| \le 1$).

\emph{(i)(a).} $\gamma_{0}$ commutes with $I\sigma_{3} = \gamma_{2}\gamma_{1}$ and anticommutes with $\kappa$. So $WA_{0}\phi = \gamma_{0}\gamma_{0}\phi\gamma_{0}(I\sigma_{3})^{-1}\gamma_{0} = \phi(I\sigma_{3})^{-1} = A_{0}W\phi$, and $WB_{k}\phi = -\kappa\phi(I\sigma_{3})^{-1}\gamma_{0}$ while $B_{k}W\phi = -\gamma_{0}\kappa\gamma_{0}\phi\gamma_{0}(I\sigma_{3})^{-1} = +\kappa\phi\gamma_{0}(I\sigma_{3})^{-1}$: $W$ commutes with the mass part and anticommutes with the momentum part, and it commutes with the right unit. For $\tilde\phi(\eta) = W\phi(-\eta)$ this gives $\tilde\phi' = -W(M(-\eta)A_{0} + B_{k})\phi(-\eta) = (-M(-\eta)A_{0} + B_{k})\tilde\phi$, which is $G(\eta)\tilde\phi$ for every solution exactly when $M(-\eta) = -M(\eta)$. $W$ is the outermorphism of the reflection $x \mapsto -\gamma_{0}x\gamma_{0}$ on even elements (checked on the six bivectors). Right multiplication by $\sigma_{1} = \gamma_{1}\gamma_{0}$ anticommutes with right multiplication by $I\sigma_{3}$, commutes with $A_{0}$ (two sign changes: $\sigma_{1}\gamma_{0} = -\gamma_{0}\sigma_{1}$, $\sigma_{1}I\sigma_{3} = -I\sigma_{3}\sigma_{1}$) and carries $B_{k}$ to $B_{-k}$, so $c$ maps solutions of momentum $\vec k$ to solutions of momentum $-\vec k$ for every $M$, antilinearly. Hence $\theta G_{k}(-\eta) = -G_{-k}(\eta)\theta$ for odd $M$, and $\theta$ is an antilinear symmetry. $W$, $c$ and the unit are orthogonal for the scalar product. Left multiplication by $I_{4}$ anticommutes with $A_{0}$ and commutes with $B_{k}$ (the conventions' statement that it flips the mass), so it is not a symmetry and plays no part. All of this is verified in $\Cl(1,3)$ on the even blades at four samples of $(\vec k, M)$.

\emph{(i)(b).} An orthogonal map $O$ of the vectors of $\Cl(3,0)$ extends to an outermorphism with $O(a\wedge b) = Oa\wedge Ob$ and $O(I) = \det(O)I$, so $O\bigl(-2I(a\wedge b)\bigr) = \det(O)\bigl(-2I(Oa\wedge Ob)\bigr)$. For $\tilde s(\eta) = Os(-\eta)$: $\tilde s' = -O\bigl(-2I(h(-\eta)\wedge s(-\eta))\bigr) = -\det(O)\bigl(-2I(Oh(-\eta)\wedge\tilde s)\bigr)$. This equals $-2I(h(\eta)\wedge\tilde s)$ for every solution exactly when $-\det(O)\,Oh(-\eta) = h(\eta)$ for all $\eta$, that is, with $h(\pm\eta) = k\sigma_{1} \pm M\sigma_{3}$, for every $k$ and, since $\mu > 0$, some $M \ne 0$, $-\det(O)O\sigma_{1} = \sigma_{1}$ and $\det(O)O\sigma_{3} = \sigma_{3}$. If $\det O = -1$: $O\sigma_{1} = \sigma_{1}$, $O\sigma_{3} = -\sigma_{3}$, and then $O\sigma_{2} = +\sigma_{2}$; this is $O_{A}$. If $\det O = +1$: $O\sigma_{1} = -\sigma_{1}$, $O\sigma_{3} = \sigma_{3}$, $O\sigma_{2} = -\sigma_{2}$; this is $O_{B}$. The check confirms that of the eight diagonal sign maps exactly these two intertwine the equation. A unitary $U$ of a sector acting on states as $\rho \mapsto U\rho U^{-1}$ gives a rotation of Bloch vectors, an antiunitary one an improper map, and the exchange $\rho \mapsto 1 - \rho$ of particles and holes gives $s \mapsto -s$. $W$ commutes with the right unit, the spatial rotors and the translations, hence preserves each sector, where it is a unitary commuting with $\sigma_{3}$ and anticommuting with $\sigma_{1}$, so a phase times $\sigma_{3}$, with $\sigma_{3}v\sigma_{3} = O_{B}v$. As an automorphism $W$ therefore induces $O_{B}$, and it sends the wall-instant vacuum $\sigma_{1}$ to the filled sea $-\sigma_{1}$: it reverses the energy. $\theta = cW$ is antiunitary on the solution space and a symmetry; since $c$ carries momentum $\vec k$ to $-\vec k$, it maps the sector $(\vec k, h)$ onto a sector $(-\vec k, h')$, and with both sectors written in the stored two-level form (the same $H$, since $H$ depends on $k$ only) it is an antiunitary map between them intertwining the equation with $\eta \mapsto -\eta$, so on Bloch vectors it is the improper one of the two, $O_{A}$, from the sector $(\vec k, h)$ to the sector $(-\vec k, h')$. $W$ followed by the exchange induces $-O_{B} = O_{A}$ within each sector. For a state of the isotropic class the two conditions coincide. Both keep $\sigma_{1}$. None of this uses $\mu > 0$ except the count of two: at $\mu = 0$ the condition is only $-\det(O)\,O\sigma_{1} = \sigma_{1}$, met by every improper orthogonal map fixing $\sigma_{1}$ and every rotation reversing it; these arise from $\theta$ and $W$ composed with the chiral rotations, left multiplication by $\exp(I_{4}\alpha)$, which commutes with $B_{k}$ and is a symmetry of the massless equation only.

\emph{(i)(c).} A point-fixing conjugation acts as $\tilde s(\eta) = O's(\eta)$ and must intertwine the equation at each $\eta$: $\det(O')O'h = h$ for all $h$ in the plane of $\sigma_{1}, \sigma_{3}$. If it reverses the energy, $O'h = -h$, so $\det O' = -1$ and $O'\sigma_{2} = -\sigma_{2}$: $O' = -1$, whose only fixed vector is $0$, the tracial state. Being antilinear and commuting with the translations it carries momentum $\vec k$ to $-\vec k$, so invariance reads $v_{-\vec k} = -v_{\vec k}$ and gives $v = 0$ in the isotropic class (without isotropy it also allows the vacuum on one half of the momenta with the filled sea on the other). Followed by the exchange it is the identity. This is the type of the stored conjugation $C = \mathrm{Ad}(I_{6})\circ\kappa$ as Premise \claimref{dark-matter.premise-cpt-covariance} describes it (it fixes every translation and is antilinear, so it reverses the energy); that description is quoted, not rederived, and the premise is not used.

\emph{(ii).} Both $v(\eta)$ and $Ov(-\eta)$ solve the same linear equation, so they agree for all $\eta$ exactly when they agree at $\eta = 0$: $Ov(0) = v(0)$. For $O_{A}$ this is $v(0)\cdot\sigma_{3} = 0$, for $O_{B}$ it is $v(0) = t\sigma_{3}$. The check verifies both on power series: with $v(0)\cdot\sigma_{3} = 0$ the components along $\sigma_{1}, \sigma_{2}$ are even in $\eta$ and the one along $\sigma_{3}$ odd, to order $24$, and not otherwise. The pure states $\pm\sigma_{3}$ have $\tfrac12(1 \mp \sigma_{3}\cdot\sigma_{1}) = \tfrac12$ over $\sigma_{1}$. The disk $v\cdot\sigma_{3} = 0$ is convex and invariant under $v \mapsto -v$.

\emph{(iii).} At $\mu = 0$, $h = k\sigma_{1}$ is constant and $v' = -2I(k\sigma_{1}\wedge v)$ rotates $v$ about $\sigma_{1}$ at rate $2k$: the stationary solutions are $v = t\sigma_{1}$, all with $v\cdot\sigma_{3} = 0$. An invariant $v(0) = a\sigma_{1} + b\sigma_{2}$ with $b \ne 0$ has $v'(0) = 2kb\,\sigma_{3} \ne 0$ (checked on the series). The conformal vacuum is $t = 1$, the filled sea $t = -1$, the tracial state $t = 0$, and a diagonal charge-symmetric deposit with occupation $n(k)$ over the conformal vacuum is $t = 1 - 2n$. An improper orthogonal map fixing $\sigma_{1}$ is the reflection along a unit vector $n$ orthogonal to $\sigma_{1}$, its invariant states are those with $v(0)\cdot n = 0$, and the statements above hold with $n$ in place of $\sigma_{3}$: every $t\sigma_{1}$ is invariant and the invariant stationary Fock states are $\pm\sigma_{1}$.

\emph{(iv)(a).} $\hat h' = k\mu\,\omega^{-2}\,(c\sigma_{3} - \sigma\sigma_{1})$ has $|\hat h'| = k\mu/\omega^{2}$, integrable over the line with integral $\pi$, and $\hat h\wedge\hat h = 0$, so $\hat h$ solves the equation up to the residual $\hat h'$. With $R(\eta, \eta')$ the rotation that the equation generates, $s_{\mathrm{out}}(\eta) = \hat h(\eta) + \int_{\eta}^{\infty}R(\eta, \eta')\hat h'(\eta')\,d\eta'$ is a solution with $|s_{\mathrm{out}} - \hat h| \le \int_{\eta}^{\infty}k\mu\,\omega^{-2} \to 0$; it is unique because the difference of two solutions has constant length, and $|s_{\mathrm{out}}| = \lim|\hat h| = 1$. For the series write $D = \omega^{-1}d/d\eta$; from $\omega' = \mu\sigma$: $Dc = -\delta c\sigma$, $D\sigma = \delta c^{2}$, $D\delta = -2\delta^{2}\sigma$, so $D$ maps polynomials in $(c, \sigma)$ times $\delta^{n}$ to polynomials times $\delta^{n+1}$, and the equation is $Ds = -2I(\hat h\wedge s)$. Put $d_{0} = \hat h$ and, given $d_{0}, \dots, d_{n}$, let $e$ be $D(\delta^{n}d_{n})/\delta^{n+1}$; if $e\cdot\hat h = 0$, the part of $d_{n+1}$ orthogonal to $\hat h$ is $\tfrac12I(\hat h\wedge e)$, the unique solution of $-2I(\hat h\wedge x) = e$ orthogonal to $\hat h$, and its part along $\hat h$ is fixed by $|S|^{2} = 1$ at order $\delta^{n+1}$. The check runs this recursion to $N = 6$ with exact polynomials modulo $c^{2} + \sigma^{2} = 1$: the solvability condition holds at each order, the norm is $1$ order by order, $d_{1} = \tfrac12c\,\sigma_{2}$, and the residual $DS_{N} + 2I(\hat h\wedge S_{N})$ of the truncated sum is exactly of order $\delta^{N+1}$. (Solvability at every order follows from the norm condition: $|S_{n}|^{2} = 1 + O(\delta^{n+1})$ gives $S_{n}\cdot DS_{n} = O(\delta^{n+2})$, whose coefficient of $\delta^{n+1}$ is $\hat h\cdot e + \sum_{j=1}^{n}d_{j}\cdot\bigl(-2I(\hat h\wedge d_{n+1-j})\bigr)$; the sum changes sign under $j \mapsto n + 1 - j$, so it vanishes and $\hat h\cdot e = 0$.) So $S_{N}' + 2I(h\wedge S_{N}) = \omega\,\delta^{N+1}P_{N}(c, \sigma)$ with $P_{N}$ a vector polynomial, bounded on the circle $c^{2} + \sigma^{2} = 1$ by the sum $C_{N}$ of the moduli of its coefficients, and $S_{N} - \hat h \to 0$ as $\eta \to \infty$. Since the flow is a rotation, $|s_{\mathrm{out}} - S_{N}|(\eta) \le \int_{\eta}^{\infty}C_{N}\mu^{N+1}\omega^{-2N-1}\,d\eta' \le C_{N}\,(\mu/k^{2})^{N}\int_{-\infty}^{\infty}(1 + x^{2})^{-N-1/2}dx$, for every real $\eta$ and $N \ge 1$. Expanding $S_{N}$ in powers of $1/k$ at fixed $\eta$ gives a formal solution with leading term $\sigma_{1}$, which by the uniqueness clause of the stored result (its clause (i)) is the stored adiabatic series through the orders determined; so $s_{\mathrm{out}}$ agrees with that series to every order, the sense in which the stored results speak of an adiabatic state of infinite order. Under $\eta \mapsto -\eta$ one has $\sigma \mapsto -\sigma$ with $c$ and $\delta$ fixed, and by (i)(b) $O_{A}S(-\eta)$ is again a formal solution of unit norm with leading term $\hat h$; the recursion determines such a solution uniquely, so $O_{A}d_{n}(c, -\sigma) = d_{n}(c, \sigma)$ for every $n$ (the check confirms it through $n = 6$: the parts along $\sigma_{1}, \sigma_{2}$ are even in $\sigma$ and the part along $\sigma_{3}$ odd), so $S_{N}(0)\cdot\sigma_{3} = 0$ and $|s_{\mathrm{out}}(0)\cdot\sigma_{3}| \le C_{N}'(\mu/k^{2})^{N}$. At the wall $S(0) = \sigma_{1} + \tfrac12\delta\sigma_{2} - \tfrac18\delta^{2}\sigma_{1} + \dots$ with $\delta = \mu/k^{2}$ (checked), so the wall-instant vacuum $\sigma_{1}$ has $\tfrac12(1 - S(0)\cdot\sigma_{1}) = \mu^{2}/(16k^{4}) + O(k^{-8})$ over $s_{\mathrm{out}}$, the stored value with $M_{1} = \mu$.

\emph{(iv)(b).} $s_{\mathrm{in}}(0) = O_{A}s_{\mathrm{out}}(0)$ differs from $s_{\mathrm{out}}(0)$ only in the sign of the $\sigma_{3}$ part $x = s_{\mathrm{out}}(0)\cdot\sigma_{3}$, so $s_{\mathrm{in}}\cdot s_{\mathrm{out}} = 1 - 2x^{2}$ and $q = x^{2}$ (constant in $\eta$ by the rotation property). Their sum is $2\sqrt{1 - q}\,\hat m$ with $\hat m$ a unit vector orthogonal to $\sigma_{3}$ when $q < 1$. A unit vector $e$ orthogonal to $\sigma_{3}$ has $e\cdot s_{\mathrm{out}}(0) = e\cdot s_{\mathrm{in}}(0) = \sqrt{1 - q}\cos\chi$, $\chi$ its angle from $\hat m$; a unit vector with equal products against $s_{\mathrm{in}}(0)$ and $s_{\mathrm{out}}(0)$ is orthogonal to their difference $2x\sigma_{3}$, so for $x \ne 0$ the converse holds. For $|v(0)| \le 1$ orthogonal to $\sigma_{3}$, $v(0)\cdot s_{\mathrm{out}}(0) = \sqrt{1 - q}\,v(0)\cdot\hat m \in [-\sqrt{1 - q}, \sqrt{1 - q}]$, with the upper end only at $v(0) = \hat m$; this gives the bounds, the profile of $t\,\hat m$ and $4n_{0}(1 - n_{0}) = 1 - (1 - q) = q$. The check verifies these identities on rational points of the sphere.

\emph{(iv)(c). The value of $q$ (classical; Weber functions).} Numbers are in $\mathrm{span}\{1, I\}$. With $H = M\sigma_{3} + k\sigma_{1}$ and components $u_{1}, u_{2}$ along the eigenvectors of $\sigma_{3}$, $Iu_{1}' = Mu_{1} + ku_{2}$, $Iu_{2}' = ku_{1} - Mu_{2}$, and eliminating $u_{1}$ gives $u_{2}'' + (k^{2} + \mu^{2}\eta^{2} - I\mu)u_{2} = 0$ (the check verifies both second-order forms on power series with Gaussian rational coefficients). With $z = e^{I\pi/4}\sqrt{2\mu}\,\eta$, $z^{2} = 2I\mu\eta^{2}$, this is Weber's equation $w'' + (p + \tfrac12 - \tfrac14z^{2})w = 0$ with $p = -1 - I\nu/2$, $\nu = k^{2}/\mu$ (checked). The Bloch vector of $u$ has $\sigma_{3}$ part $(|u_{1}|^{2} - |u_{2}|^{2})/(|u_{1}|^{2} + |u_{2}|^{2})$, and $\hat h \to \pm\sigma_{3}$ as $\eta \to \pm\infty$, so $s_{\mathrm{out}}$ is the solution with $u_{2}/u_{1} \to 0$ at $+\infty$, and, evaluating the constant $q = \tfrac12(1 - s_{\mathrm{in}}\cdot s_{\mathrm{out}})$ at $\eta \to -\infty$ where $s_{\mathrm{in}} \to -\sigma_{3}$, $q = \lim_{\eta\to-\infty}|u_{1}|^{2}/(|u_{1}|^{2} + |u_{2}|^{2})$ on that solution. Take $u_{2} = D_{p}(z)$. For $|\arg z| < \tfrac34\pi$, $D_{p}(z) \sim z^{p}e^{-z^{2}/4}$; on $\arg z = \tfrac14\pi$ ($\eta \to +\infty$), $|e^{-z^{2}/4}| = 1$ and $|z^{p}| = |z|^{-1}e^{\pi\nu/8} \to 0$, so this is the out solution; there $Iu_{2}' = Mu_{2}(1 + O(\eta^{-2}))$, so $ku_{1} = Iu_{2}' + Mu_{2}$ gives $|u_{1}| \to (2\mu\eta/k)\,e^{\pi\nu/8}/(\sqrt{2\mu}\,\eta)$ and the conserved norm is $|u_{1}|^{2} + |u_{2}|^{2} = (2/\nu)e^{\pi\nu/4}$. For $\eta \to -\infty$, $\arg z = -\tfrac34\pi$, the connection formula $D_{p}(z) = e^{-I\pi p}D_{p}(-z) + \sqrt{2\pi}\,\Gamma(-p)^{-1}e^{-I\pi(p+1)/2}D_{-p-1}(Iz)$ applies with $\arg(-z) = \tfrac14\pi$ and $\arg(Iz) = -\tfrac14\pi$: the first term has modulus $|z|^{-1}e^{\pi\nu/8} \to 0$, and the second, with $-p - 1 = I\nu/2$, $|(Iz)^{I\nu/2}| = e^{\pi\nu/8}$, $|e^{-I\pi(p+1)/2}| = e^{-\pi\nu/4}$ and $|\Gamma(1 + I\nu/2)|^{2} = \tfrac12\pi\nu/\sinh(\tfrac12\pi\nu)$ (the reflection formula), has squared modulus $(4/\nu)\sinh(\tfrac12\pi\nu)\,e^{-\pi\nu/4}$. Dividing by the norm, $|u_{2}|^{2}/(|u_{1}|^{2} + |u_{2}|^{2}) \to 2\sinh(\tfrac12\pi\nu)e^{-\pi\nu/2} = 1 - e^{-\pi\nu}$, so $q = e^{-\pi\nu} = e^{-\pi k^{2}/\mu}$. The check verifies the exponent bookkeeping ($-\tfrac12 + \tfrac14 - \tfrac14 = -\tfrac12$ in units of $\pi\nu$) and the identity $2\sinh(y)e^{-y} = 1 - e^{-2y}$ on rational values of $e^{-y}$; the asymptotic form, the connection formula (whose case $p = -1$ is $D_{-1}(z) + D_{-1}(-z) = \sqrt{2\pi}e^{z^{2}/4}$) and the reflection formula are classical (Whittaker and Watson, A Course of Modern Analysis, sections 16.5 to 16.52; cited from memory). This is the Landau-Zener transition probability (Zener 1932, from memory). A numerical integration of the Bloch equation from $\eta = 60$ to the wall (computed in floating point, not part of the check; started on $\hat h + \tfrac12\delta c\,\sigma_{2}$ and integrated with a Runge-Kutta method of order eight at relative tolerance $10^{-11}$) reproduces $q = e^{-\pi\nu}$ and $n_{0}$ to nine digits at five values of $(k, \mu)$, and (computed, by numerical quadrature) $\int_{0}^{\infty}x^{2}\bigl(1 - \sqrt{1 - e^{-x^{2}}}\bigr)dx/(2\pi^{2}) = 0.012760$; since $n_{0}$ is one half of the bracket, the floor carries the comoving number density $0.006380\,(\mu/\pi)^{3/2}$ per sector and $0.012760\,(\mu/\pi)^{3/2}$ summed over the two helicities (substitute $k = x\sqrt{\mu/\pi}$ in $\int d^{3}k\,n_{0}/(2\pi)^{3}$).

\emph{(v)(a).} On a collar where $a > 0$ is smooth, the difference of the two-point functions of two states of the class, at $(\eta, \vec x)$ and $(\eta', \vec x')$, is $a(\eta)^{-3/2}a(\eta')^{-3/2}$ times an integral over $\vec k$ and sum over sectors of $\exp(-I\sigma_{3}\vec k\cdot(\vec x - \vec x'))$ against the sector operator $U_{k}(\eta, 0)\,\tfrac12(v_{k}(0) - r_{k}(0))\,U_{k}(\eta', 0)^{-1}$, carried into the field's components by an isometry depending on the direction of $\vec k$ only. $U_{k}$ is unitary and $\partial_{\eta}U_{k} = -IHU_{k}$ with $|H| = \omega \le k + |M|$, so a derivative of total order $j$ in $(\eta, \eta', \vec x, \vec x')$ multiplies the integrand by a factor bounded by $C_{j}(1 + k)^{j}$ uniformly for $\eta, \eta'$ in a compact set (the derivatives of $H$ are those of $M$, bounded there). With $|v_{k}(0) - r_{k}(0)| = |v_{k} - r_{k}|$ constant in $\eta$ and $\int k^{N}|v_{k} - r_{k}|\,d^{3}k < \infty$ for every $N$, every derivative of the integral converges absolutely and uniformly on compact sets, so the kernel is smooth. For Fock states $|v - r|^{2} = 2(1 - v\cdot r) = 4n$, and by the Cauchy-Schwarz inequality $\int_{k > 1}k^{N}\sqrt{n}\,d^{3}k \le (\int k^{2N+4}n)^{1/2}(\int_{k>1}k^{-4})^{1/2}$.

\emph{(v)(b).} $s_{\mathrm{out}}$ agrees with the adiabatic series to every order by (iv)(a) and is a homogeneous Fock state of the class, so by Fact \claimref{synthesis.imported-dirac-hadamard-form-propagates-and-adiabatic-states-are-hadamard}(c) it is Hadamard on $\eta > \varepsilon$, where $a$ is smooth and positive (the identification of that agreement with the Fact's adiabatic states of infinite order is the one the stored massive result makes). A kernel differing from a Hadamard one by a smooth kernel has the same wave front set, so is Hadamard (Fact \claimref{synthesis.imported-microlocal-hadamard-condition}). For $v = t(k)e_{\chi(k)}$: $|v - s_{\mathrm{out}}| \le (1 - t) + |e_{\chi} - s_{\mathrm{out}}|$ and $|e_{\chi} - s_{\mathrm{out}}|^{2} = 2 - 2\sqrt{1 - q}\cos\chi \le 2(1 - \cos\chi) + 2q$, with $q \le C_{N}'^{2}(\mu/k^{2})^{2N}$ by (iv)(a) (or $q = e^{-\pi k^{2}/\mu}$ by (iv)(c)); so (v)(a) applies when $1 - t$ and $1 - \cos\chi$ decrease faster than every power. The wall-instant vacuum is the instantaneous vacuum at $\eta = 0$, where $M_{1} = \mu \ne 0$, and is not Hadamard by clause (ii) of the stored massive result; the tracial state is not Hadamard by the stored moment criterion (Theorem \claimref{synthesis.deposit-over-the-conformal-vacuum-is-hadamard-exactly-when-every-occupation-moment-is-finite}) in the massless case. For $\mu > 0$ this is argued, not proved: (v)(a) is a sufficient condition only and no necessity criterion for massive fields is stored; the argument is that the tracial state's two-point function is one half of the anticommutator distribution in both orderings, whose wave front set contains both $C^{(+)}$ and $C^{(-)}$, so neither ordering has the one-sided wave front set. The same holds, argued, for the Fock states of $O_{B}$, whose occupation over $s_{\mathrm{out}}$ tends to one half. That (CW) does not imply (HD) rests on the wall-instant vacuum, by the stored clause, which itself is given Fact (c); that (HD) does not imply (CW) rests on $s_{\mathrm{out}}$, Hadamard given the Fact and not invariant for $\mu > 0$. The Fermi-Dirac profile at $\mu = 0$ is Hadamard by clause (iii) of that criterion, and it is a state $t\sigma_{1}$ of (iii).

\emph{(vi).} Among invariant states the occupation over $s_{\mathrm{out}}$ is at least $n_{0}$, with equality only at $v(0) = \hat m$ (iv)(b), a Fock state; at $\mu = 0$, $s_{\mathrm{in}} = s_{\mathrm{out}} = \sigma_{1}$, $q = 0$, $\hat m = \sigma_{1}$. Without the invariance the minimum over all states is $0$, at $s_{\mathrm{out}}$, which is invariant only if $q = 0$.

\emph{(vii).} By (i)(b), $\theta$ acts from the sector $(\vec k, h)$ to the sector $(-\vec k, h')$ as $s(\eta) \mapsto O_{A}s(-\eta)$, so by the argument of (ii) a state with one vector per sector is invariant exactly when $v_{-\vec k, h'}(0) = O_{A}v_{\vec k, h}(0)$ for each such pair (applied twice it returns $v_{\vec k, h}$, as $O_{A}^{2} = 1$): one vector of each pair is free. $s_{\mathrm{out}}$ and $s_{\mathrm{in}}$ depend on $k$ only. With $v = v_{\vec k, h}(0)$, $n_{\vec k, h} = \tfrac12(1 - v\cdot s_{\mathrm{out}}(0))$ and $n_{-\vec k, h'} = \tfrac12(1 - O_{A}v\cdot s_{\mathrm{out}}(0)) = \tfrac12(1 - v\cdot s_{\mathrm{in}}(0))$, since $O_{A}$ is a symmetric orthogonal map and $O_{A}s_{\mathrm{out}}(0) = s_{\mathrm{in}}(0)$. Their sum is $1 - \tfrac12v\cdot(s_{\mathrm{in}}(0) + s_{\mathrm{out}}(0)) = 1 - \sqrt{1 - q}\,v\cdot\hat m$, at least $1 - \sqrt{1 - q} = 2n_{0}$ because $|v| \le 1$, with equality only for $v = \hat m$, where $O_{A}v = v$ and both sectors carry $s_{0}$. The choice $v = s_{\mathrm{out}}(0)$ gives the pair $(s_{\mathrm{out}}, s_{\mathrm{in}})$, with occupations $0$ and $\tfrac12(1 - s_{\mathrm{in}}\cdot s_{\mathrm{out}}) = q$, and $0 < n_{0}$ for $q > 0$. The check verifies the sum formula, the bound and this example on rational points. $W$ preserves each sector and so does the exchange of particles and holes, so for that conjugation the arguments of (ii), (iv)(b) and (vi) apply to each sector's vector separately.

\emph{Reading (argued, not part of the statement).} What (CW) selects and what it does not. With the orientation datum and the minimum of (vi) it gives a vacuum part with no free number: the conformal vacuum for massless fields and the profile $n_{0}(k^{2}/ma')$ for massive ones. It gives no temperature, no thermal bath and no spectrum: every charge-symmetric diagonal deposit of a massless field is invariant, the thermal ones at every temperature included, so the hot deposit is left exactly as free as (HD) leaves it. This agrees with the two obstructions recorded on op.deposit-state-is-hadamard that it meets. (O1): the invariant set is symmetric under $v \mapsto -v$, the conjugation does not tell the vacuum from the filled sea, the only state fixed by every conjugation of (i) is the tracial one, and the point-fixing energy-reversing type fixes nothing else; lifting the deposit off one half takes the orientation datum, here the choice of $s_{\mathrm{out}}$ over $-s_{\mathrm{out}}$ in (vi). That datum is the one (N$_{+}$) assumes and (HD) already carries where it reaches (\claimref{note.synthesis.n-plus-orientation-datum}), so the bundle of (CW), purity, least occupation and the orientation does not source it: it trades (HD) for the same datum and three further clauses. (O5): the invariance adds no condition to stationarity for massless fields, so a hot deposit is neither selected nor excluded. The mirror. If the mirror is read as the far side, it is not a second algebra: the fields at $-\eta$ are the Cauchy evolution of the fields at $\eta$, so hypothesis (B) of Theorem \claimref{synthesis.cold-mirror-needs-non-equilibrium-descent} (two decoupled algebras with a product dynamics) does not hold for the pair, and by (iv)(b) every invariant state reads the same occupations on the two sides (over $s_{\mathrm{in}}$ at $-\eta$ and over $s_{\mathrm{out}}$ at $\eta$; for a stationary massless state the same $t(k)$), a temperature ratio $1$. If instead the mirror is a second field species exchanged with this one by the conjugation, an invariant quasi-free state gives the two species conjugate marginals, again equal thermal profiles for stationary massless states, and the invariance restricts neither marginal by itself. In both readings a mirror colder than this tower in a quasi-free state is not (CW)-invariant; the exchange-invariant state with unequal temperatures is the mixture of clause (iii) of that theorem, not quasi-free, with components that are not invariant. The series' cold mirror (temperature ratio below $0.448$) and the deposit law of the mirror's temperature therefore require that (CW) is not adopted for the thermal part of the deposit, or that the realised state is one component of an invariant mixture; this is a tension with the series' modelling of the mirror, not a contradiction with a stored theorem, since the theorem's hypothesis (B) fails or is untouched. It is a cost of (CW) alone: no stored premise asks the state to be invariant under the exchange, P2 being scoped to sector labels and not to thermal states within a sector (\claimref{dark-matter.p2-indifference}), and the conjugation premise being a covariance of structures and actions. The background. (RD) and (DF) prescribe the gravitational field, with $a = a'\eta$ exactly through $\eta = 0$; with the expected stress of a Hadamard state of the content as the source the exact radiation field does not reach the wall (the result filed alongside, synthesis.trace-anomaly-leaves-no-exact-radiation-field-at-a-tame-wall, given its hypotheses (A) and (E)), so the present result is for a test field on a prescribed background, and the reflection $W$, which needs $M$ exactly odd at $\eta = 0$, is not available on a background corrected in that way. Transport-absent: the content is the standard theory of the Dirac field with a mass linear in conformal time. Prior art: L. Boyle, K. Finn and N. Turok, Phys. Rev. Lett. 121 (2018) 251301, cited from memory, and their longer paper arXiv:1803.08930, whose locators were confirmed in review against the arXiv text: there CPT invariance of the vacuum on the doubled radiation field leaves a free function of the momentum for spinors (their Eq. (3.48), which pairs the mode of momentum $p$ with that of $-p$), the occupation over the late-time vacuum is $(1 - \cos 2\eta\cos\lambda)/2$ in their notation (Eq. (3.49)), the function is fixed by minimising the particle number and the energy, and with $|\beta| = e^{-\pi p^{2}/2\gamma}$ between the early-time and late-time vacua (Eq. (5.13)) the number density summed over the two helicities is $0.01276\,(\gamma/\pi)^{3/2}$ (Eqs. (5.15), (5.16)); clauses (iv)(b), (iv)(c) and (vi) restate this in the store's form. What is added here is the classification (i) against the stored conjugation, the mixed states and the floor (ii) and (iv)(b), the pair bound (vii), the elementary bound (iv)(a), and the Hadamard statements (v) with the wall-instant vacuum as an invariant state that is not Hadamard.
